\documentclass[10pt,a4paper]{amsart}
\usepackage{amsmath,amssymb,amsthm,graphicx}
\usepackage{verbatim}
\newenvironment{reference}{\comment}{\endcomment}
%\newenvironment{reference}{}{}
\newenvironment{slogan}{\comment}{\endcomment}
\newenvironment{history}{\comment}{\endcomment}

% For commutative diagrams we use Xy-pic
\usepackage[all]{xy}

% We use 2cell for 2-commutative diagrams.
\xyoption{2cell}
% Theorem environments.
%
\theoremstyle{plain}
\newtheorem{theorem}[subsection]{Theorem}
\newtheorem{proposition}[subsection]{Proposition}
\newtheorem{lemma}[subsection]{Lemma}

\theoremstyle{definition}
\newtheorem{definition}[subsection]{Definition}
\newtheorem{example}[subsection]{Example}
\newtheorem{exercise}[subsection]{Exercise}
\newtheorem{situation}[subsection]{Situation}

\theoremstyle{remark}
\newtheorem{remark}[subsection]{Remark}
\newtheorem{remarks}[subsection]{Remarks}

\numberwithin{equation}{subsection}

% Macros
%
\def\lim{\mathop{\mathrm{lim}}\nolimits}
\def\colim{\mathop{\mathrm{colim}}\nolimits}
\def\Spec{\mathop{\mathrm{Spec}}}
\def\Hom{\mathop{\mathrm{Hom}}\nolimits}
\def\Ext{\mathop{\mathrm{Ext}}\nolimits}
\def\SheafHom{\mathop{\mathcal{H}\!\mathit{om}}\nolimits}
\def\SheafExt{\mathop{\mathcal{E}\!\mathit{xt}}\nolimits}
\def\Sch{\mathit{Sch}}
\def\Mor{\mathop{\mathrm{Mor}}\nolimits}
\def\Ob{\mathop{\mathrm{Ob}}\nolimits}
\def\Sh{\mathop{\mathit{Sh}}\nolimits}
\def\NL{\mathop{N\!L}\nolimits}
\def\CH{\mathop{\mathrm{CH}}\nolimits}
\def\proetale{{pro\text{-}\acute{e}tale}}
\def\etale{{\acute{e}tale}}
\def\QCoh{\mathit{QCoh}}
\def\Ker{\mathop{\mathrm{Ker}}}
\def\Im{\mathop{\mathrm{Im}}}
\def\Coker{\mathop{\mathrm{Coker}}}
\def\Coim{\mathop{\mathrm{Coim}}}

% Boxtimes
%
\DeclareMathSymbol{\boxtimes}{\mathbin}{AMSa}{"02}

%
% Macros for moduli stacks/spaces
%
\def\QCohstack{\mathcal{QC}\!\mathit{oh}}
\def\Cohstack{\mathcal{C}\!\mathit{oh}}
\def\Spacesstack{\mathcal{S}\!\mathit{paces}}
\def\Quotfunctor{\mathrm{Quot}}
\def\Hilbfunctor{\mathrm{Hilb}}
\def\Curvesstack{\mathcal{C}\!\mathit{urves}}
\def\Polarizedstack{\mathcal{P}\!\mathit{olarized}}
\def\Complexesstack{\mathcal{C}\!\mathit{omplexes}}
% \Pic is the operator that assigns to X its picard group, usage \Pic(X)
% \Picardstack_{X/B} denotes the Picard stack of X over B
% \Picardfunctor_{X/B} denotes the Picard functor of X over B
\def\Pic{\mathop{\mathrm{Pic}}\nolimits}
\def\Picardstack{\mathcal{P}\!\mathit{ic}}
\def\Picardfunctor{\mathrm{Pic}}
\def\Deformationcategory{\mathcal{D}\!\mathit{ef}}
\title{Nilpotence of an element in a residue fibre is detected by the characteristic polynomial}
\author{The Stacks Project Authors}
\date{stacks-project:algebra.tex:local-snapshot-20261012; GNU FDL 1.2 or later}
\begin{document}
\maketitle
\noindent\textit{Editorial note.} All delivered mathematical text is the author's own, copied from the pinned Stacks Project chapter file algebra.tex; no statement or proof step was rewritten, reordered or removed.\par\medskip
\noindent\textit{Editorial note.} Preamble provenance, as required for this source: the frozen snapshot stages the chapter's own preamble file beside it (preamble.tex, 217 lines, 8116 bytes, the same shared Stacks-wide preamble the published chow and etale-cohomology cards use); the first line of the pinned chapter file is \textbackslash{}input\{preamble\}, so no preamble text lives inside the chapter itself. The card takes that file's lines 8-18 (the verbatim package and the three \textbackslash{}comment-based environments reference/slogan/history at lines 9-12, section 13.2, together with the Xy-pic setup) and its lines 154 to the runtime end of the file (the native theorem declarations and mathematical macros). Not one character was changed, no package was added, no macro was defined, and the excerpt serves only as preamble, never as a span coordinate.\par\medskip
\noindent\textit{Editorial note.} The same-chapter pointers are located at their \textbackslash{}label lines inside the pinned chapter file; the cross-chapter pointers are located in the project's own tag list staged in the snapshot (sources/stacks-chow/tags.txt), at the verbatim lines carrying their canonical tags.\par\medskip
\noindent\textit{Editorial note.} The counted spans contain no \textbackslash{}ref and no \textbackslash{}cite, so the card carries no navigation stub and no bibliography.\par\medskip
\noindent\textit{Allowed standard prerequisites (editorial).}\begin{itemize}
\item Cayley-Hamilton for the multiplication map. Used in the counted proof (lines 5786-5800). Ties nilpotence of the multiplication map to the vanishing of the characteristic polynomial.
\item Behaviour of the characteristic polynomial under base change to a residue field. Used in the counted proof (lines 5787-5808). Identifies the characteristic polynomial of the fibre with the reduction of the original one.
\end{itemize}\medskip
\begin{lemma}
\label{lemma-characteristic-polynomial-prime}
Let $R \to A$ be a ring homomorphism.
Assume $A \cong R^{\oplus n}$ as an $R$-module.
Let $f \in A$. The multiplication map $m_f: A
\to A$ is $R$-linear and hence
has a characteristic polynomial
$P(T) = T^n + r_{n-1}T^{n-1} + \ldots + r_0 \in R[T]$.
For any prime
$\mathfrak{p} \in \Spec(R)$, $f$ acts nilpotently on $A
\otimes_R \kappa(\mathfrak{p})$ if and only if $\mathfrak p \in
V(r_0, \ldots, r_{n-1})$.
\end{lemma}

\begin{proof}
This follows quite easily once we prove that the characteristic
polynomial $\bar P(T) \in \kappa(\mathfrak p)[T]$ of the
multiplication map $m_{\bar f}: A \otimes_R \kappa(\mathfrak p) \to
A \otimes_R \kappa(\mathfrak p)$ which multiplies elements of $A
\otimes_R \kappa(\mathfrak p)$ by $\bar f$, the image of $f$ viewed in
$\kappa(\mathfrak p)$, is just the image of $P(T)$ in
$\kappa(\mathfrak p)[T]$. Let $(a_{ij})$ be the matrix of the map
$m_f$ with entries in $R$, using a basis $e_1, \ldots, e_n$
of $A$ as an $R$-module.
Then, $A \otimes_R \kappa(\mathfrak p) \cong (R \otimes_R
\kappa(\mathfrak p))^{\oplus n} = \kappa(\mathfrak p)^n$, which is
an $n$-dimensional vector space over $\kappa(\mathfrak p)$ with
basis $e_1 \otimes 1, \ldots, e_n \otimes 1$. The image $\bar f = f
\otimes 1$, and so the multiplication map $m_{\bar f}$ has matrix
$(a_{ij} \otimes 1)$. Thus, the characteristic polynomial is
precisely the image of $P(T)$.

\medskip\noindent
From linear algebra, we know that a linear transformation acts
nilpotently on an $n$-dimensional vector space if and only if the
characteristic polynomial is $T^n$ (since the characteristic
polynomial divides some power of the minimal polynomial). Hence,
$f$ acts nilpotently on $A \otimes_R \kappa(\mathfrak p)$ if and
only if $\bar P(T) = T^n$. This occurs if and only if $r_i \in
\mathfrak p$ for all $0 \leq i \leq n - 1$, that is when $\mathfrak p \in
V(r_0, \ldots, r_{n - 1}).$
\end{proof}

\end{document}
